\section{Related Work}
\label{sec:related}

The problem studied here draws on three lines of research, namely subgroup discovery, the clustering of longitudinal trajectories, and the interpretation of clusters through rules. Each line supplies part of what is needed, and the gap that remains is the reason this work combines them.

\paragraph{Subgroup discovery and exceptional model mining.} Subgroup discovery searches for subsets of a population that stand out with respect to a property of interest, and it returns each subset as an interpretable description rather than as an opaque score \citep{atzmueller2015subgroup,herrera2011overview}. The idea goes back to the early pattern-mining systems of \citet{klosgen1996explora} and \citet{wrobel1997midos}, and efficient exhaustive search was later made practical \citep{atzmueller2006sdmap}. Contrast-set mining, emerging-pattern mining, and subgroup discovery were unified as supervised descriptive rule discovery by \citet{novak2009supervised}. Exceptional model mining generalises the idea from a single target attribute to a model fitted on several targets, so that a subgroup is exceptional when a model learned on it differs from the model learned on the whole population \citep{leman2008emm,duivesteijn2016emm}. Recent work casts the search as continuous optimisation, as in Syflow, which maximises the divergence between the subgroup and the population and learns interpretable rules end to end \citep{xu2024syflow}. Other extensions reach sequential and temporal targets, including subgroups with exceptional transition behaviour modelled by Markov chains \citep{lemmerich2016transition} and local exceptionality detection in time series \citep{knobbe2008local}. These methods produce readable descriptions, yet their targets are scalars, distributions, or transition models rather than a shape traced by a continuous quantity over a shared calendar, which is the target considered here.

\paragraph{Clustering of longitudinal trajectories.} A large body of work groups entities by the shape of a quantity that they trace over time. Group-based trajectory modelling assumes that a population is a mixture of a few latent trajectory classes and estimates that mixture \citep{nagin2005group,nagin2010group}, and growth mixture modelling extends the idea with random effects inside each class \citep{muthen1999finite}. Distance-based clustering of longitudinal data has been packaged for applied use, for example the longitudinal $k$-means of \citet{genolini2010kml} and the broader framework of \citet{denteuling2023latrend}. The wider time-series clustering literature offers many shape-based distances, including dynamic time warping \citep{berndt1994dtw}, and it has been reviewed at length \citep{aghabozorgi2015timeseries}. What these methods return is cluster membership. An entity is told which class it belongs to, but not a rule that would let a reader recognise the members of that class from their attributes, which is the missing half that a descriptive method must supply.

\paragraph{Interpreting clusters with rules.} Recent work makes clusters interpretable by describing them with axis-aligned trees. \citet{dasgupta2020explainable} give an explainable $k$-means and $k$-medians procedure in which a small tree recovers a near-optimal clustering, and a growing literature fits decision trees to explain the partitions of standard clustering algorithms, surveyed by \citet{hongfei2024interpretableclustering}. This turns a clustering into rules after the fact. SEAL is close in spirit, since it also explains a clustering with a shallow tree, but it searches over many clusterings and keeps the one whose tree explanation is both faithful and short, so the description enters the objective rather than following as a post-hoc report. The broader case for preferring inherently interpretable models over explanations of opaque ones is argued by \citet{rudin2019stop} and surveyed by \citet{molnar2022interpretable}.

\paragraph{Position of this work.} SEAL keeps the descriptive rule of subgroup discovery and the temporal target of trajectory clustering, and it produces the grouping and its rule together rather than in two stages. In keeping with the descriptive-mining tradition, where discovered patterns are judged by their interpretability and validity, we demonstrate the method on case studies in which the rules can be checked against documented domain structure. Two elements of the descriptor construction receive particular attention. The first is the engineering of temporal features so that a rule can speak about how a covariate moved rather than only its final value. The second is a non-circularity discipline that keeps a rule from restating the definition of its own target. Both matter most in the household health case study, where the target is a ratio of two quantities that could otherwise leak into the descriptors.
